\chapter{Support Projections and Central Support}

This chapter records the next formalization target after Sakai's Theorem 1.9.1: completion of
Section 1.10 through Proposition 1.10.7.  Propositions 1.10.1 and 1.10.2 are already formalized as
Proposition \ref{prop:left_ideals_in_wstar_proj_Sak_1_10_1}, Corollary
\ref{cor:right_ideals_in_wstar_proj_Sak_1_10_1}, and Proposition
\ref{prop:proj_compl_lat_wstar_Sak_1_10_2}.  The remaining results construct left, right, and
self-adjoint support projections, identify ultraweakly closed two-sided ideals with central
projections, and prove the orthogonality theorem for central supports.

This is also the necessary bridge to Sakai's spectral resolution in Section 1.11, where the
spectral family is defined by
\[
  e(\lambda)=s\bigl((\lambda 1-h)^{+}\bigr).
\]

The local Mathlib version and the original LeanOA repository have no operator-algebra support
projection or central-support construction to reuse.  The implementation should nevertheless be
mostly connective tissue around existing general objects:
\texttt{LinearMap.ker}, Mathlib's native left ideals \texttt{Ideal}, native
\texttt{TwoSidedIdeal}, \texttt{Set.center}, the complete lattice of projections already proved in
this repository, and Mathlib's nonunital star-subalgebra closure API.  In particular, annihilators
must be introduced at the ring-theoretic level, not as sets specialized to normed spaces, and a
two-sided ideal must not be encoded by a new repository-specific structure.

\section{Algebraic annihilators}

The topology used later should be absent from the basic definitions.  For a subset $S$ of a
semiring $R$, write
\[
  \operatorname{Ann}_{\ell}(S)=\{x\in R\mid xs=0\text{ for every }s\in S\},\qquad
  \operatorname{Ann}_{r}(S)=\{x\in R\mid sx=0\text{ for every }s\in S\}.
\]
The first is a left ideal of $R$, while the second is represented by a left ideal of
$R^{\mathrm{op}}$, consistently with Mathlib's convention.

\begin{definition}
  \label{def:one_sided_annihilators}
  \notready
  Let $R$ be a semiring and $S\subseteq R$.  Define the left annihilator of $S$ as an
  \texttt{Ideal R} and the right annihilator as an \texttt{Ideal (MulOpposite R)}.  For a singleton,
  these are the kernels of right multiplication by $a$ in $R$ and by
  $\operatorname{op}(a)$ in $R^{\mathrm{op}}$, respectively.
\end{definition}

These definitions should be built from infima of \texttt{LinearMap.ker}; their carrier and
singleton simp lemmas should reduce to the displayed pointwise conditions.  This formulation
requires no star operation, order, norm, topology, or scalar field.

\begin{definition}
  \label{def:two_sided_annihilators}
  \notready
  Let $R$ be a ring.  The annihilator of a left ideal $I$,
  \[
    \{x\in R\mid xs=0\text{ for every }s\in I\},
  \]
  and the annihilator of a right ideal $J$,
  \[
    \{x\in R\mid sx=0\text{ for every }s\in J\},
  \]
  are native Mathlib \texttt{TwoSidedIdeal R}s.
\end{definition}

\begin{proof}
  \notready
  Absorption on one side follows by associativity and zero multiplication.  On the other side,
  use that $I$ is closed under left multiplication and that $J$, encoded in $R^{\mathrm{op}}$, is
  closed under right multiplication.  Package the resulting carriers with
  \texttt{TwoSidedIdeal.mk'} and expose membership simp lemmas.
\end{proof}

\begin{lemma}
  \label{lem:annihilators_ultraweak_closed}
  \notready
  \uses{def:one_sided_annihilators,def:two_sided_annihilators,lem:star_left_mul_right_mul_sig_cts_Sak_1_7_8}
  Let $M$ be a $C^{*}$--algebra with specified complete predual $P$.  The left and right
  annihilators of an arbitrary subset of $M$ are ultraweakly closed.  The two-sided annihilator of
  any left or right ideal is also ultraweakly closed when viewed through its underlying complex
  submodule.
\end{lemma}

\begin{proof}
  \notready
  Each annihilator is an intersection of kernels of maps of the form $x\mapsto xa$ or
  $x\mapsto ax$.  Fixed left and right multiplication are already ultraweakly continuous, so the
  kernels are closed and arbitrary intersections remain closed.  The proof should use the
  transported submodule API and should not repeat a net calculation.
\end{proof}

\begin{lemma}
  \label{lem:ultraweak_op_homeomorph}
  \notready
  The canonical maps $\operatorname{op}:M\to M^{\mathrm{op}}$ and
  $\operatorname{unop}:M^{\mathrm{op}}\to M$ induce mutually inverse complex-linear
  homeomorphisms for the ultraweak topologies determined by a specified predual $P$.  Consequently
  a subset is ultraweakly closed if and only if its image in the opposite algebra is ultraweakly
  closed.
\end{lemma}

\begin{proof}
  \notready
  The specified predual on $M^{\mathrm{op}}$ is already transported from the predual on $M$ along
  Mathlib's linear isometric equivalence \texttt{op}/\texttt{unop}.  Apply the functorial API for
  weak topologies to that equivalence.  This should be a general transport lemma for specified
  preduals, not a fact tied to ideals.
\end{proof}

\section{Left, right, and self-adjoint support}

The public support projections should not retain a choice of predual.  For a $W^{*}$--algebra
$M$, define them using the canonical complete lattice
$\{p:M\mid p\text{ is a projection}\}$:
\[
  l(a)=\bigwedge\{p\mid pa=a\},\qquad
  r(a)=\bigwedge\{p\mid ap=a\}.
\]
The annihilator argument supplies the substantive fact that these infima themselves fix $a$.
This separates the choice-free public objects from the explicit predual used in their proofs.

\begin{definition}
  \label{def:left_right_support_Sak_1_10_3} (Sakai 1.10.3)
  \notready
  \uses{prop:proj_compl_lat_wstar_Sak_1_10_2}
  For $a$ in a $W^{*}$--algebra $M$, its \emph{left support} $l(a)$ is the infimum of the
  projections $p$ satisfying $pa=a$, and its \emph{right support} $r(a)$ is the infimum of the
  projections $q$ satisfying $aq=a$.  Both constructions take values in the existing subtype of
  star projections.
\end{definition}

\begin{proposition}
  \label{prop:left_right_support_characterization}
  \notready
  \uses{def:left_right_support_Sak_1_10_3,lem:annihilators_ultraweak_closed,prop:left_ideals_in_wstar_proj_Sak_1_10_1,cor:right_ideals_in_wstar_proj_Sak_1_10_1}
  For every $a\in M$,
  \[
    l(a)a=a,\qquad ar(a)=a.
  \]
  Moreover, for every projection $p$,
  \[
    l(a)\leq p\ \Longleftrightarrow\ pa=a,
    \qquad
    r(a)\leq p\ \Longleftrightarrow\ ap=a.
  \]
  Equivalently, the singleton annihilators satisfy
  \[
    \operatorname{Ann}_{\ell}(\{a\})=M(1-l(a)),\qquad
    \operatorname{Ann}_{r}(\{a\})=(1-r(a))M.
  \]
\end{proposition}

\begin{proof}
  \notready
  By Lemma \ref{lem:annihilators_ultraweak_closed} and Proposition
  \ref{prop:left_ideals_in_wstar_proj_Sak_1_10_1}, write the left annihilator as $Me$ for a unique
  projection $e$.  Then $(1-e)a=a$.  If $pa=a$, the projection $1-p$ annihilates $a$, hence
  $1-p\leq e$ and $1-e\leq p$.  Thus $1-e$ has the defining universal property of $l(a)$.
  The right-support statement is the same argument in the opposite algebra.  Record both the
  universal properties and the annihilator identities as the primary API, so later proofs do not
  unfold choices of ideal generators.
\end{proof}

\begin{lemma}
  \label{lem:support_star_compatibility}
  \notready
  \uses{prop:left_right_support_characterization}
  For every $a\in M$,
  \[
    l(a^{*})=r(a),\qquad r(a^{*})=l(a).
  \]
  If $h=h^{*}$, these projections agree.  Their common value is denoted $s(h)$ and is called the
  support projection of $h$.
\end{lemma}

\begin{proof}
  \notready
  Taking adjoints converts $pa=a$ into $a^{*}p=a^{*}$ because $p$ is self-adjoint.  Apply the two
  order characterizations in Proposition \ref{prop:left_right_support_characterization} and use
  antisymmetry.  This part is algebraic once the support universal properties are available.
\end{proof}

Several useful consequences should be included with this node: supports of zero and one,
invariance under multiplication by a nonzero scalar, and monotonicity of support on nonnegative
elements.  They should be proved from the universal properties or continuous functional calculus,
as appropriate, rather than from the implementation of the infimum.

\section{The nonunital ultraweak algebra generated by an element}

Sakai's ``$W^{*}$--subalgebra generated by $h$'' in Proposition 1.10.4 need not contain the
ambient identity.  Its internal identity is precisely $s(h)$.  Therefore the existing unital
ultraweak-closure construction is not the correct formal object; the parallel construction for
Mathlib's nonunital star subalgebras is required.

\begin{lemma}
  \label{lem:internal_unit_is_projection}
  \notready
  Let $N$ be a star-stable nonunital subalgebra of a star ring.  If $N$ has an internal identity,
  then the ambient value of that identity is a star projection.
\end{lemma}

\begin{proof}
  \notready
  Idempotence is the unit law.  Apply the right unit law to the adjoint of the identity and take
  adjoints to prove self-adjointness.  The current proof of this fact is private inside the
  ultraweak left-ideal argument; it should be promoted to a reusable algebraic lemma and that file
  refactored to use it.
\end{proof}

\begin{definition}
  \label{def:nonunital_ultraweak_closure}
  \notready
  Let $S$ be a nonunital star subalgebra of a $C^{*}$--algebra with specified predual.  Define its
  ultraweak transport, the predicate that it is ultraweakly closed, its ultraweak closure, and the
  nonunital ultraweak star algebra generated by a set.
\end{definition}

\begin{proposition}
  \label{prop:nonunital_ultraweak_closure_api}
  \notready
  \uses{def:nonunital_ultraweak_closure,lem:wstar_unital,lem:internal_unit_is_projection}
  The nonunital ultraweak closure contains the original algebra, is ultraweakly closed, and is
  minimal among ultraweakly closed nonunital star subalgebras containing it.  It inherits a
  nonunital $C^{*}$--algebra structure; if the predual is complete, it is internally unital by the
  existing closed-submodule unitality theorem, and the ambient value of that identity is a star
  projection.  Ultraweak closure preserves commutativity.
\end{proposition}

\begin{proof}
  \notready
  Mirror the established unital ultraweak-closure API, reusing Mathlib's topological closure for
  nonunital star subalgebras.  The proof of internal unitality is an application of Lemma
  \ref{lem:wstar_unital} to the underlying closed submodule and its tautological linear isometry,
  not a second proof of Sakai's unitality theorem.
\end{proof}

\begin{proposition}
  \label{prop:support_mem_generated_wstar_Sak_1_10_4} (Sakai 1.10.4)
  \notready
  \uses{lem:support_star_compatibility,prop:left_right_support_characterization,prop:nonunital_ultraweak_closure_api}
  Let $h=h^{*}$ and let $C$ be the nonunital ultraweak star algebra generated by $h$.  Then
  $s(h)\in C$.  More strongly, when $C$ is equipped with its inherited nonunital
  $C^{*}$--algebra structure, the ambient value of its internal identity is $s(h)$.
\end{proposition}

\begin{proof}
  \notready
  Let $p\in C$ be the internal identity.  Since $ph=h$, minimality gives $s(h)\leq p$.
  The element $p-s(h)$ annihilates $h$, hence it annihilates every nonunital star polynomial in
  $h$.  Fixed multiplication is ultraweakly continuous, so it also annihilates the ultraweak
  closure $C$, and in particular $(p-s(h))p=0$.  As $s(h)\leq p$, this forces $p=s(h)$.
\end{proof}

\section{Closed two-sided ideals and central projections}

Use Mathlib's bundled \texttt{TwoSidedIdeal M} throughout this section.  Its underlying left and
right ideals are \texttt{TwoSidedIdeal.asIdeal} and
\texttt{TwoSidedIdeal.asIdealOpposite}; ultraweak closedness is imposed on their common underlying
complex submodule.

\begin{proposition}
  \label{prop:closed_two_sided_ideal_central_projection_Sak_1_10_5} (Sakai 1.10.5)
  \notready
  \uses{prop:left_ideals_in_wstar_proj_Sak_1_10_1,cor:right_ideals_in_wstar_proj_Sak_1_10_1,lem:ultraweak_op_homeomorph}
  Every ultraweakly closed two-sided ideal $J$ of a $W^{*}$--algebra $M$ has the form
  \[
    J=Mz=zM
  \]
  for a unique central projection $z$.  Formally, $J$ is a native
  \texttt{TwoSidedIdeal M}, and the conclusion identifies both its underlying left ideal with
  \texttt{Ideal.span \{z\}} and its opposite underlying left ideal with
  \texttt{Ideal.span \{MulOpposite.op z\}}.
\end{proposition}

\begin{proof}
  \notready
  Apply the left-ideal case of Proposition 1.10.1.  Transport closedness through Lemma
  \ref{lem:ultraweak_op_homeomorph} and apply the right-ideal case to obtain projections
  $z_{\ell}$ and $z_{r}$.  Since both are identities on the same ideal,
  $z_{\ell}=z_{r}=z$.  For $x\in M$, the elements $xz$ and $zx$ lie in $J$; using $z$ as the right
  and left identity of $J$ gives $xz=zxz=zx$.  Thus $z$ is central.  Uniqueness follows from
  uniqueness in either one-sided representation.
\end{proof}

\begin{corollary}
  \label{cor:central_projection_closed_ideal_order_iso}
  \notready
  \uses{prop:closed_two_sided_ideal_central_projection_Sak_1_10_5,prop:proj_compl_lat_wstar_Sak_1_10_2}
  Central projections in $M$, ordered by the ambient projection order, are order-isomorphic to
  ultraweakly closed two-sided ideals via $z\mapsto Mz$.  Consequently the subtype of central
  projections is a complete lattice.
\end{corollary}

\begin{proof}
  \notready
  Restrict the existing order isomorphism between projections and ultraweakly closed left ideals.
  Proposition \ref{prop:closed_two_sided_ideal_central_projection_Sak_1_10_5} identifies exactly
  which left ideals are two-sided.  Arbitrary intersections of closed two-sided ideals are again
  closed and two-sided, so the same infimum construction used for Proposition 1.10.2 transports a
  complete lattice structure.  The implementation should use a subtype of Mathlib's projection
  subtype together with membership in \texttt{Set.center}, rather than introduce a second notion
  of projection or centrality.
\end{proof}

\begin{definition}
  \label{def:central_support_Sak_1_10_6} (Sakai 1.10.6)
  \notready
  \uses{cor:central_projection_closed_ideal_order_iso}
  If $p$ is a projection in $M$, its \emph{central support} $c(p)$ is the infimum, in the complete
  lattice of central projections, of all central projections $z$ satisfying $p\leq z$.
\end{definition}

\begin{proposition}
  \label{prop:central_support_characterization}
  \notready
  \uses{def:central_support_Sak_1_10_6}
  The central support $c(p)$ is central, $p\leq c(p)$, and for every central projection $z$,
  \[
    c(p)\leq z\ \Longleftrightarrow\ p\leq z.
  \]
  The operation $c$ is monotone and idempotent, and $c(p)=p$ precisely when $p$ is central.
\end{proposition}

\begin{proof}
  \notready
  These are the standard infimum and closure-operator laws in the complete lattice of central
  projections.  They should be proved order-theoretically from Corollary
  \ref{cor:central_projection_closed_ideal_order_iso}; no new analytic argument is needed.
\end{proof}

\section{Orthogonality of central supports}

\begin{lemma}
  \label{lem:projection_mul_eq_zero_iff_le_complement}
  \notready
  \uses{lem:proj_sub_pos_iff_comm_eq_self}
  If $p$ and $q$ are projections in a unital $C^{*}$--algebra, then
  \[
    pq=0\ \Longleftrightarrow\ p\leq 1-q.
  \]
  Equivalently, $p$ and $q$ are disjoint in the projection order if and only if their product is
  zero.
\end{lemma}

\begin{proof}
  \notready
  Apply the existing characterization $p\leq r\iff pr=p$ to the projection $r=1-q$ and expand
  $p(1-q)$.  This lemma belongs with the general $C^{*}$--algebra projection API and must not be
  specialized to $W^{*}$--algebras.
\end{proof}

\begin{proposition}
  \label{prop:central_support_orthogonal_Sak_1_10_7} (Sakai 1.10.7)
  \notready
  \uses{def:two_sided_annihilators,lem:annihilators_ultraweak_closed,prop:closed_two_sided_ideal_central_projection_Sak_1_10_5,prop:central_support_characterization,lem:projection_mul_eq_zero_iff_le_complement}
  Let $p,q$ be projections in a $W^{*}$--algebra $M$.  If
  \[
    pMq=\{0\},\qquad\text{that is,}\qquad pxq=0\quad\text{for every }x\in M,
  \]
  then
  \[
    c(p)c(q)=0.
  \]
\end{proposition}

\begin{proof}
  \notready
  Let $J$ be the two-sided annihilator of the right ideal $pM$.  It is ultraweakly closed by
  Lemma \ref{lem:annihilators_ultraweak_closed}, and the hypothesis says $q\in J$.  By Proposition
  \ref{prop:closed_two_sided_ideal_central_projection_Sak_1_10_5}, write $J=Mz$ for a central
  projection $z$.  Then $q\leq z$, hence $c(q)\leq z$, so $pc(q)=0$.  Lemma
  \ref{lem:projection_mul_eq_zero_iff_le_complement} gives $p\leq1-c(q)$.  The complement is
  central, and minimality of $c(p)$ yields $c(p)\leq1-c(q)$, which is equivalent to
  $c(p)c(q)=0$.
\end{proof}

\begin{corollary}
  \label{cor:central_support_orthogonal_iff}
  \notready
  \uses{prop:central_support_orthogonal_Sak_1_10_7,prop:central_support_characterization}
  For projections $p,q$ in $M$,
  \[
    c(p)c(q)=0\ \Longleftrightarrow\ (\forall x\in M,\;pxq=0).
  \]
\end{corollary}

\begin{proof}
  \notready
  The reverse implication is Proposition
  \ref{prop:central_support_orthogonal_Sak_1_10_7}.  Conversely, use $p\leq c(p)$ and
  $q\leq c(q)$ together with centrality to insert $c(p)$ and $c(q)$ into $pxq$; the product then
  vanishes because $c(p)c(q)=0$.  This equivalence is the reusable API form of Sakai's
  proposition.
\end{proof}
